% Options for packages loaded elsewhere
\PassOptionsToPackage{unicode}{hyperref}
\PassOptionsToPackage{hyphens}{url}
\documentclass[12pt,a4paper]{article}
\usepackage{xcolor}
\usepackage{amsmath,amssymb}
\setcounter{secnumdepth}{3}
\usepackage{iftex}
\ifPDFTeX
\usepackage[T1]{fontenc}
\usepackage[utf8]{inputenc}
\usepackage{textcomp}
\else
\usepackage{unicode-math}
\defaultfontfeatures{Scale=MatchLowercase}
\defaultfontfeatures[\rmfamily]{Ligatures=TeX,Scale=1}
\fi
\usepackage{lmodern}
\IfFileExists{upquote.sty}{\usepackage{upquote}}{}
\IfFileExists{microtype.sty}{%
\usepackage[stretch=10,shrink=10]{microtype}
\UseMicrotypeSet[protrusion]{basicmath}
}{}
\makeatletter
\@ifundefined{KOMAClassName}{%
\IfFileExists{parskip.sty}{%
\usepackage{parskip}
}{%
\setlength{\parindent}{0pt}
\setlength{\parskip}{6pt plus 2pt minus 1pt}}
}{%
\KOMAoptions{parskip=half}}
\makeatother
\usepackage{longtable,booktabs,array}
\usepackage{multirow}
\newcounter{none}
\usepackage{calc}
\usepackage{etoolbox}
\makeatletter
\patchcmd\longtable{\par}{\if@noskipsec\mbox{}\fi\par}{}{}
\makeatother
\IfFileExists{footnotehyper.sty}{\usepackage{footnotehyper}}{\usepackage{footnote}}
\makesavenoteenv{longtable}
\usepackage{graphicx}

% ===== IMAGES: center + fit to page width =====
\setkeys{Gin}{width=0.85\textwidth,height=!,keepaspectratio}
\let\oldincludegraphics\includegraphics
\renewcommand{\includegraphics}[2][]{%
  \begin{center}\oldincludegraphics[#1]{#2}\end{center}}
% ==============================================

\makeatletter
\newsavebox\pandoc@box
\newcommand*\pandocbounded[1]{%
\sbox\pandoc@box{#1}%
\Gscale@div\@tempa{\textheight}{\dimexpr\ht\pandoc@box+\dp\pandoc@box\relax}%
\Gscale@div\@tempb{\linewidth}{\wd\pandoc@box}%
\ifdim\@tempb\p@<\@tempa\p@\let\@tempa\@tempb\fi%
\ifdim\@tempa\p@<\p@\scalebox{\@tempa}{\usebox\pandoc@box}%
\else\usebox{\pandoc@box}%
\fi%
}
\def\fps@figure{htbp}
\makeatother
\usepackage{float}
\usepackage{hyperref}
\hypersetup{hidelinks}

\begin{document}
\sloppy  % Allows slightly more flexible line breaking to avoid overfull hboxes

%% ============ CUSTOM COVER PAGE ============
\begin{titlepage}
  \centering
  \vspace*{0.4cm}

  % University logo
  \oldincludegraphics[width=0.22\textwidth]{media/image1.png}

  \vspace{0.3cm}
  {\large\bfseries University of Messina\par}
  {\normalsize Master Degree in Data Science\par}
  {\normalsize High Performance Computing --- Project Report\par}

  \vspace{0.4cm}
  \rule{\linewidth}{0.4pt}
  \vspace{0.3cm}

  {\Large\bfseries Hybrid MPI + OpenMP Monte Carlo Forecasting
  of Household Electric Power Consumption\par}

  \vspace{0.3cm}
  \rule{\linewidth}{0.4pt}

  \vspace{0.6cm}
  \begin{minipage}[t]{0.45\textwidth}
    \raggedright
    {\large\bfseries {Professor}\par}
    {\large\bfseries Prof.\ Salvatore Distefano\par}
  \end{minipage}
  \hfill
  \begin{minipage}[t]{0.45\textwidth}
    \raggedleft
    {\large\bfseries {Student}\par}
    {\large\bfseries Nehil Jain\par}
    {\large\bfseries Matricola: 562482\par}
  \end{minipage}

  \vspace{0.6cm}
  {\normalsize Academic Year 2025/2026\par}
\end{titlepage}
%% ==========================================

%% Page numbering: cover is unnumbered, ToC is roman, content restarts at 1
\pagenumbering{roman}
\tableofcontents
\newpage
\pagenumbering{arabic}

\section{Introduction}\label{introduction}
\subsection{General Introduction}\label{general-introduction}

Forecasting electric power consumption is a central problem in modern
energy management. Utilities, grid operators and households all benefit
from reliable short-term predictions of demand: they enable better
generation scheduling, demand-response programs, tariff planning and
anomaly detection. Consumption time series, however, are inherently
noisy and uncertain, driven by human behaviour, weather and appliance
usage patterns. Deterministic point forecasts are therefore of limited
value; what decision makers need is a probabilistic forecast, i.e. an
estimate of the expected consumption together with a quantification of
its uncertainty.

Monte Carlo simulation is a natural tool for this purpose. By simulating
a stochastic model of the consumption process many thousands of times,
one obtains an empirical distribution of future trajectories from which
means and confidence intervals can be derived. The accuracy of these
estimates grows with the number of simulated trajectories, which makes
the method computationally expensive and, at the same time, an excellent
candidate for parallelization.

\subsection{Introduction to the Scenario and Related Issues}

The scenario considered in this project is the 30-day probabilistic
forecast of seven electrical quantities measured in a single household:
global active power, global reactive power, voltage, global intensity
and three sub-metering energy channels. The input is the UCI
``Individual household electric power consumption'' dataset [{1}],
containing about 2 million one-minute measurements collected over almost
four years. For each of the seven quantities, the forecast must provide
the expected daily value and a 95\% confidence interval for each of the
next 30 days, at one-minute simulation resolution (1,440 steps per day).

The computational issue is the sheer size of the simulation workload.
With 50,000 Monte Carlo trajectories per quantity, 7 quantities, 30 days
and 1,440 steps per day, the program must execute approximately 15
billion stochastic update steps, each involving the generation of a
Gaussian random number and the evaluation of an exponential demand-shape
function. On a single core of a modern laptop this requires more than
eleven minutes (675.1 s measured); embedded in any interactive or
repeated-analysis workflow, such latency is unacceptable.

\subsection{Motivations}

The primary motivation of the project is to reduce the execution time of
the fore-casting computation by exploiting all the parallel resources of
modern hardware, following the two complementary paradigms studied in
the course: message pass-ing between processes (MPI) for inter-node,
distributed-memory parallelism, and multithreading within a process
(OpenMP) for intra-node, shared-memory parallelism.

The problem offers a particularly clean case study because it
naturally exposes both forms of parallelism at the same time: the 50,000
simulations are independent data items that can be partitioned across
processes (data parallelism), while the seven per-column forecasting
pipelines are independent functional units that can be executed
concurrently by threads (task parallelism).

A secondary, educational motivation is to quantify and explain the
limits of each paradigm on real hardware: the granularity ceiling of
task parallelism (only seven tasks exist), the load imbalance among
tasks of unequal duration, and the heterogeneity of a modern hybrid CPU
(performance and efficiency cores), all of which leave measurable
signatures in the experimental results of Chapter 7.

\subsection{Existing Solutions and Proposed Solution}

Monte Carlo methods are classified as ``embarrassingly parallel'' in the
literature, since trajectories are mutually independent [5].
Standard parallel formulations partition the trajectory set across
workers and combine partial statistics at the end; this approach is used
both in distributed-memory implementations based on MPI and in
shared-memory implementations based on OpenMP or GPU computing [4].
Special care is required for the random number streams: each trajectory
must receive an independent, reproducible stream regardless of which
worker executes it, otherwise the parallel result differs from the
sequential one or, worse, trajectories become statistically correlated.

The solution proposed in this project is a hybrid MPI + OpenMP
implementation that combines both decompositions in a single executable,
following the structure required by the project specification: MPI
implements data parallelism by block-partitioning the simulation index
space across ranks, while OpenMP implements task parallelism by
executing the seven column pipelines as explicit tasks with declared
dependencies inside each rank. A deterministic, position-based seeding
scheme guarantees that the sequential, pure OpenMP, pure MPI and hybrid
versions produce bit-identical statistical results, providing a built-in
correctness verification. Four versions of the program were implemented
and benchmarked: sequential (baseline $T_1$), pure OpenMP, pure MPI, and hybrid.

\section{Problem Rationale}

\subsection{Description of the Problem}

Given the historical measurements of the seven electrical quantities,
the program must produce, for each quantity c \{0,\ldots,6\} and each
future day d \{1,\ldots,30\}, the expected daily average value and a
95\% confidence interval. The consumption process is modelled as a
discrete-time stochastic process with drift, volatility and a
deterministic intra-day demand pattern. Each Monte Carlo trajectory
simulates 30 days at one-minute resolution; the per-day averages of all
trajectories are then aggregated into the required statistics. The
problem is computation-bound: the input parameters are small (three scalars per column), the output is
small (three $7\times30$ matrices), but the simulation kernel performs on the
order of 15 billion floating-point updates.

\subsection{Current Sequential Solution and Algorithm}

The sequential algorithm proceeds in four phases. (1) Loading: the CSV
dataset is parsed and each column is normalized with a z-score
transformation, which makes the model parameters comparable across
columns of different physical units. (2) Parameter estimation: for each
column, the starting value, the volatility and the drift are estimated
from the most recent window of the historical series. (3) Simulation:
for each column, 50,000 independent trajectories are generated; every
trajectory evolves the value minute by minute with the update val $\leftarrow$val +
drift/1440 + z\,$\cdot$\,vol\,$\cdot$\,pattern(hour)\,$\cdot$\,0.01, where z is a standard
Gaussian sample produced by a Box-Muller transform over a linear
congruential generator, and pattern(hour) is a deterministic shape with
morning (08:00) and evening (18:00) demand peaks. Per-day sums and
squared sums are accumulated. (4) Reporting: sums are converted into
means and 95\% confidence intervals, de-normalized back to physical
units and written to the report. The structure of the sequential program
is shown in Figure 1; the highlighted simulation block accounts for
almost the entire execution time (T= 675.1 s) and is the target of
parallelization.

\includegraphics[width=4.05278in,height=6.62292in]{media/image2.png}

\begin{quote}
\emph{Figure 1: Flowchart of the sequential algorithm. After loading and
normalizing the dataset, the seven columns are processed one at a time;
for each column, parameters are estimated, 50,000 Monte Carlo
trajectories are simulated (30 days$\times$ 1,440 one-minute steps), and
per-day statistics are computed. The highlighted simulation block
dominates the execution time and is the target of parallelization.}
\end{quote}

\subsection{Existing Parallel Solutions, State of the Art and Related Work}

Parallel Monte Carlo is a classical topic in high performance computing.
Rosenthal [5] formalizes the conditions under which Monte Carlo
estimators can be computed in parallel without bias, the key requirement
being independent random streams per work unit. Foster's
methodology [6] --- partitioning, communication, agglomeration and
mapping --- provides the standard design framework that this report
follows chapter by chapter. In distributed-memory settings, the dominant
pattern is the partition of trajectories across MPI ranks followed by a
collective reduction of partial sums [4]; in shared-memory settings,
OpenMP worksharing or task constructs distribute trajectories or
higher-level functional units across threads. Hybrid MPI+OpenMP designs
combine the two, using MPI between nodes and OpenMP within a node, which
reduces the number of MPI ranks, the memory footprint and the
communication volume compared with a pure MPI configuration of equal
worker count. The present project follows this hybrid pattern, with one
specific design twist dictated by the assignment: the MPI level carries
the data decomposition while the OpenMP level carries the task
decomposition.

Monte Carlo estimation rests on the law of large numbers: if $X_1$,
\ldots{} , $X_n$ are i.i.d.\ random variables with finite expectation, then
in this project X is the daily average consumption of a
column on a given future day, and each realization is produced by one
simulated trajectory. Independence across trajectories is also what
makes the method parallel-friendly: any subset of trajectories can be
computed by any processing element, and partial sums simply add up. The
random number generator used is a 32-bit linear congruential generator
combined with the Box-Muller transform; while statistically simple, it
is fully adequate for this workload and, crucially, it is seedable per
trajectory, which the parallel design exploits.

\subsection{Adopted Stochastic Model}

Each column c is described by three parameters estimated from the
historical data: a starting level, a volatility and a linear drift. A
trajectory evolves the normalized value for 43,200 one-minute steps (30
days). The intra-day modulation $\text{pattern}(h) = 1 +
0.3\cdot\exp(-0.5\cdot((h-8)/2)^2) + 0.4\cdot\exp(-0.5\cdot((h-18)/2)^2)$ amplifies
volatility around the morning and evening demand peaks observed in the
dataset. For each day the trajectory contributes its daily average to
two accumulators, sum\_d and sum2\_d, from which the mean and the 95\%
confidence interval (mean $\pm$ 1.96\,$\cdot$\,$\sigma$) are obtained at the end. The seed of
trajectory s of column c is the pure function seed = c\,$\cdot$\,NUM\_SIMS + s +
1: it does not depend on which process or thread executes the
trajectory, which (i) makes every parallel version reproduce the sequential results exactly and (ii)
assigns disjoint random streams to different columns, avoiding
cross-column correlation.

\subsection{Decomposition of the Problem}

The problem exposes two orthogonal axes of decomposition, and the
project specification assigns one technology to each. The first axis is
the simulation index space: the 50,000 trajectories of every column are
independent data items of identical computational cost, so the index
range [0, NUM\_SIMS) can be partitioned into contiguous blocks, one
per MPI rank --- this is the data decomposition. Rank r of P processes
the global indices [r\,$\cdot$\,N/P, (r+1)\,$\cdot$\,N/P) for all seven columns and
accumulates partial sums. Because the seed is a function of the global
index only, the union of the blocks reproduces the sequential
computation exactly, and the partial sums of the blocks add up to the
full-range sums.

The second axis is the set of columns: the seven per-column pipelines
are functionally independent units that read and write disjoint data ---
this is the task decomposition, implemented with OpenMP explicit tasks
inside each rank. In the pure OpenMP version each column pipeline is
itself split into three dependent tasks --- parameter estimation,
simulation and statistics --- chained by OpenMP depend clauses so that
the runtime enforces their order automatically; in the hybrid version,
where the parameters are estimated once and broadcast by MPI, the
in-rank pipeline reduces to two dependent tasks (simulation and local
reduce), chained by the same kind of depend clause. Figure 2 shows the
parallel structure obtained by combining the two decompositions; it is
the parallel counterpart of the sequential flowchart of Figure 1, with
the column loop replaced by concurrent column pipelines and the
trajectory range replaced by per-rank blocks.

\includegraphics[width=4.06458in,height=5.95972in]{media/image3.png}

\begin{quote}
\emph{Figure 2: Parallelized algorithm (hybrid version). Blue regions
mark the MPI level: the 50,000 simulations are block-decomposed across P
ranks (data parallelism), with one initial broadcast of the model
parameters and one final sum-reduction of the partial statistics. Red
regions mark the OpenMP level inside each rank: the seven column
pipelines execute as independent tasks with explicit dependency clauses
(task parallelism); in the hybrid each pipeline is a two-stage simulate
$\rightarrow$ local-reduce chain (the estimation stage is performed once on rank 0
and broadcast), joined by a taskwait before the reduction. The pure MPI
version corresponds to this figure with one thread per rank; the pure
OpenMP version corresponds to a single rank executing the full
simulation range.}
\end{quote}

\subsection{Partitioning Strategy}

At the MPI level a static block partitioning is used: rank r receives
the contiguous block [r\,$\cdot$\,N/P, (r+1)\,$\cdot$\,N/P) of trajectory indices,
computed with integer arithmetic so that block sizes differ by at most
one trajectory. Since every trajectory performs exactly the same number
of operations, static block partitioning yields near-perfect load
balance at any rank count and requires no scheduling communication. At
the OpenMP level the partitioning unit is the column task: seven tasks
are created by a single thread and executed by the team under the
runtime's task scheduler. The granularity of this level
is fixed by the problem (seven columns), which bounds the useful thread
count at seven; this granularity limit is analysed quantitatively in
Chapter 7.

\subsection{Rationale for Combined Data and Task Parallelism}

The two decompositions are complementary in exactly the dimensions where
each one is weak. The task decomposition has a hard ceiling of seven
workers and suffers from load imbalance because the seven tasks have
unequal durations (Section 5.2); but within a shared-memory node it
requires no communication at all, since tasks write disjoint slots of
shared arrays. The data decomposition scales to any number of workers
with near-perfect balance, but each additional MPI rank adds a process,
with its own address space and its participation in the collective
operations. The MPI/hybrid composition uses few MPI ranks (data
parallelism between them) and several threads inside each rank (task
parallelism), obtaining the scalability of the former and the
zero-communication sharing of the latter. The experimental results
confirm the rationale: on this 8-core machine the data-parallel MPI
level scales best (fastest at 8 ranks), while the hybrid 4$\times$2
configuration is the strongest hybrid and remains competitive; on a
larger cluster, where pure MPI would need many more ranks, the hybrid
composition is expected to lead.

\subsection{Challenges in Decomposition and Partitioning}

Three challenges emerged. First, reproducibility: a naive parallel
design that seeds trajectories by thread or rank identifier produces
results that depend on the worker count; the position-based seeding
scheme described in Section 2.5 was introduced precisely to decouple the
random streams from the execution layout. Second, the aggregation
semantics: means and confidence intervals cannot be averaged across
partial results, so the design must exchange raw sums and finalize the
statistics only once, after the global reduction (Chapter 4). Third, the
fixed task granularity: with only seven tasks of unequal duration, the
task level alone cannot use more than seven workers efficiently; the
chosen remedy is architectural --- placing the scalable data
decomposition at the other level --- rather than an artificial
refactoring of the task structure.

\section{Communication}

\subsection{Inter-Process Communication (MPI)}

The MPI communication pattern of the program is minimal and fully
collective; no point-to-point messages are used. Two events occur per
run. At the beginning, rank 0 --- the only rank that reads the dataset
--- broadcasts the estimated model parameters to all ranks with
MPI\_Bcast: 21 doubles of parameters (start, volatility, drift for each
of the seven columns) plus 14 doubles of normalization constants, 280
bytes in total. At the end, the partial per-day sums of all ranks are
combined with a single MPI\_Reduce with operation MPI\_SUM towards rank
0: two arrays of $7\times30$ doubles (sum\_d and sum2\_d), 3.36 KB in total.
The interaction pattern is therefore a star centred on rank 0, shown in
Figure 4: worker ranks never communicate with each other, and no
communication at all takes place during the simulation phase.

\includegraphics[width=4.05278in,height=2.79583in]{media/image4.png}

\begin{quote}
\emph{Figure 4: Task interaction graph for the MPI level.} \emph{Nodes
represent MPI processes; edges represent actual data exchanges, in
contrast with Figure 3 where edges represent precedence constraints. The
pattern is a star centred on rank 0: one MPI\_Bcast of 35 doubles at the
start and one MPI\_Reduce (SUM) of 420 doubles at the end. The total
communication volume is constant (3.6 KB) regardless of the number of
simulations, so the communication-to-computation ratio tends to zero as
the problem size grows.}
\end{quote}

\subsection{Intra-Process Communication (OpenMP)}

Inside a rank, the column-pipeline tasks communicate only implicitly,
through the shared address space. The design eliminates synchronization
needs by construction: each task of column c reads and writes only the
slots of index c in
the shared per-column arrays (params[c], sums[c], sums2[c],
and the result slots), so the three stages of a column hand data to each
other through these slots while different columns never share memory.
Since no two tasks ever touch the same memory locations, no critical
sections, atomics or reduction clauses are required during execution.
The only synchronization points are the task-creation single region, the
depend clauses that order the three stages of each column, and the final
taskwait, whose role is precedence enforcement rather than data
exchange. This zero-contention layout is the main efficiency advantage
of the task level and one of the reasons for combining it with MPI
instead of using MPI alone.

\subsection{Task Dependencies}

The precedence structure of the computation is captured by the task
dependency graph of Figure 3. A first node T\_L (dataset loading and
normalization) precedes everything; each of the seven columns is then a
three-stage pipeline --- estimate\_params $\rightarrow$ run\_column\_chunk
$\rightarrow$ finalize\_stats --- and the seven pipelines lie on parallel paths with
no edges between them, because each pipeline reads and writes only its
own column's data; a final node T\_R (aggregation and
report) requires the completion of all seven pipelines. In the
implementation the precedence edges are written explicitly in the source
code through OpenMP depend clauses on per-column tokens: the estimation
task declares depend(out: params[c]), the simulation task declares
depend(in: params[c]) and depend(out: sums[c]), and the
statistics task declares depend(in: sums[c]); a final taskwait
implements the fan-in before the report (or, in the MPI versions, before
the reduction). Two quantitative properties follow directly from the
shape of the graph: the maximum degree of concurrency is 7 (seven
independent column pipelines can run at once), which upper-bounds the
useful thread count of the task level; and the critical path has
constant length --- the three stages of one column (T\_L $\rightarrow$ estimate
$\rightarrow$ simulate $\rightarrow$ finalize $\rightarrow$ T\_R) --- independent of the number of workers. In
the hybrid version the estimation stage is hoisted out of the task graph
(it is performed once on rank 0 and broadcast), so each in-rank column
pipeline is the two-stage chain simulate $\rightarrow$ local-reduce, declared with
depend(out: sums[c]) followed by depend(in: sums[c]); the
maximum degree of concurrency (seven) and the per-column independence
are unchanged.

\includegraphics[width=4.05833in,height=2.78194in]{media/image5.png}

\begin{quote}
\emph{Figure 3: Task dependency graph. Nodes are tasks and directed
edges express precedence constraints (``must complete before''). Each
column is a three-stage pipeline (estimate $\rightarrow$ simulate $\rightarrow$ finalize); the
seven column pipelines lie on parallel paths with no edges between them,
so the maximum degree of concurrency is 7, while the critical path is
the three stages of a single column. The fan-out and fan-in edges
correspond directly to the OpenMP depend(in)/depend(out) clauses on the
per-column tokens and the final taskwait in the source code. In the
hybrid version a reduced form of this graph is instantiated inside every
MPI rank: since the parameters are broadcast once by rank 0, the
estimation stage is dropped and each column becomes a two-stage pipeline
(simulate $\rightarrow$ local reduce) operating on the rank's simulation block,
while the seven pipelines remain mutually independent.}
\end{quote}

\subsection{Balancing Communication and Computation}

The communication-to-computation ratio of the design is asymptotically
ideal. The communication volume is constant --- two collective
operations totalling about 3.6 KB --- while the computation grows
linearly with NUM\_SIMS: with the benchmark workload, the program
executes 15 billion update steps against a few microseconds of message
transfer. Measured wall times confirm the imbalance in favour of
computation: in every parallel run the difference between the simulation
time and the total wall time (which includes both collectives) is at
most a few milliseconds out of hundreds of seconds. As a consequence,
the speedup limits observed in Chapter 7 are not caused by
communication, but by hardware-level effects (frequency scaling and core
heterogeneity) and by the task-granularity ceiling.

\subsection{Challenges in Communication}

The main communication-related challenge was semantic rather than
volumetric: choosing what to reduce. Reducing finished statistics (means
or confidence bounds) is incorrect, because the mean of partial means with equal
weights only coincides with the global mean if all blocks have exactly
equal size, and confidence bounds do not compose at all. The
implementation therefore reduces the raw accumulators sum\_d and
sum2\_d, which are exactly additive, and computes the statistics only
once on rank 0 after the reduction. A second, minor challenge is
thread-safety of the MPI library: the hybrid version initializes MPI
with MPI\_Init\_thread requesting MPI\_THREAD\_FUNNELED and confines
every MPI call to the master thread outside the parallel regions, which
satisfies the funnelled contract on any compliant implementation.

\section{Aggregation}

\subsection{Role of Aggregation in the Algorithm}

Aggregation transforms the contributions of 50,000 independent
trajectories, possibly computed by different threads and processes, into
the final per-day statistics of each column. It operates at three nested
levels: within a trajectory (1,440 per-minute values are averaged into
one daily value), within a worker (daily values of all local
trajectories are accumulated into sum\_d and sum2\_d), and across
workers (partial accumulators are sum-reduced into global accumulators).

\subsection{Aggregation Strategy and Implementation}

The strategy follows a single invariant: only raw additive quantities
cross worker boundaries. Each trajectory contributes its daily average a
to sum\_d[day] and $a^2$ to sum2\_d[day]; these arrays are exactly
additive across any partition of the trajectory set, so the MPI level
combines them with one MPI\_Reduce(MPI\_SUM). Rank 0 then applies the
finalization once: mean = sum\_d/N, variance = sum2\_d/N $-$ mean$^2$, and the
95\% interval mean $\pm$ 1.96\,$\cdot$\,$\sigma$. At the OpenMP level no reduction operation
is needed at all, because tasks own disjoint columns and write disjoint
accumulator slots; aggregation across threads is achieved purely by
memory layout. The mean-of-means pitfall --- finalizing statistics on
partial data and averaging them afterwards --- is thereby avoided by
construction.

\subsection{Frequency and Granularity of Aggregation}

Aggregation frequency is the minimum possible: exactly once per run, at
the end of the simulation phase. The granularity is coarse on the
communication side (one message of 420 doubles per rank) and fine on the
computation side (per-day accumulation inside the innermost loop). This
choice maximizes the arithmetic intensity between synchronization
points; the alternative of periodic partial reductions would only add
latency without any benefit, since no worker needs intermediate global
results.

\subsection{Challenges in Aggregation}

Two details required care. First, numerical robustness: the variance
formula sum2/N $-$ mean$^2$ can produce small negative values due to
floating-point cancellation when the variance is close to zero (this
actually happens for the nearly constant Sub\_metering\_1 column); the
implementation clamps negative variances to zero before taking the
square root. Second, floating-point associativity: parallel summation
changes the order of additions, which in general can alter the last bits
of the result. In this design the per-column accumulation order within a
block is identical to the sequential order, and block sums are combined
in rank order by the reduction; experimentally, all configurations
reproduced the sequential statistics to every printed digit (Section
6.6).

\section{Mapping}

\subsection{Mapping Strategy for Processes and Threads}

The mapping follows the standard hybrid scheme. MPI ranks are mapped by
the runtime (Microsoft MPI's mpiexec) to the cores of
the machine, one process per worker slot; inside each rank,
omp\_set\_num\_threads fixes the team size and the OpenMP runtime maps
threads to hardware threads. Column tasks are mapped to threads
dynamically by the OpenMP task scheduler: any idle thread of the team
picks the next ready task. With seven tasks and t threads, the schedule
consists of 7/t waves; dynamic assignment lets a thread that finishes a
short task immediately steal a remaining one, which partially
compensates the duration imbalance among columns.

\subsection{Load Balancing}

The two levels have opposite balance characteristics, and the measured
data quantifies both. At the MPI level, blocks of equal size imply equal
work, so the data level is balanced by construction. At the OpenMP level
the seven tasks are far from equal: the sequential per-column times
range from 62.3 s (Sub\_metering\_2) to 148.2 s (Global\_active\_power),
a 2.38$\times$ spread, shown in Figure 8. Since the algorithm performs an
identical number of arithmetic operations per column, the spread is a
hardware effect: columns whose normalized values evolve close to zero
generate denormalized floating-point intermediates, which the CPU
processes one to two orders of magnitude slower than normalized numbers.
The consequence for scheduling is direct: the task-parallel makespan can
never be smaller than the longest task, so the task level is
upper-bounded by $T_1$ / $\max_c(t_c)$ = 675.1 / 148.2 $\approx$ 4.56, well below the
ideal 7.

\includegraphics[width=4.05833in,height=2.16042in]{media/image6.png}

\begin{quote}
\emph{Figure 8: Per-column task duration measured in the sequential run.
Although every column executes the same number of operations, durations
range from 62.3 s to 148.2 s (2.38$\times$), mainly because near-zero
normalized values produce denormalized floating-point intermediates that
the CPU processes much more slowly. The longest task lower-bounds the
task-parallel makespan.}
\end{quote}

\subsection{Impact of Mapping on Performance}

The test machine's CPU is heterogeneous --- 4
performance cores plus 4 efficiency cores --- and the mapping of
processes to core types proved to be a first-order performance factor.
Three measured effects demonstrate it. (i) With 8 MPI ranks, four ranks
necessarily run on efficiency cores; although 8 ranks (186.0 s) still
improve on 4 ranks (250.3 s) because the extra workers add real
throughput, the efficiency falls to 45\% --- the efficiency-core ranks
and the lower all-core turbo frequency prevent the speedup from
doubling. (ii) With only 2 ranks, the operating system scheduler does
not always keep both processes on performance cores; repeated runs
measured around 354 s, above the 330 s expected from two clean
performance cores, indicating occasional placement of a rank on an
efficiency core. (iii) Under full thread concurrency the per-task
durations inflate uniformly --- with 7 OpenMP threads every task took
250 s against 62--148 s in isolation --- because the all-core turbo
frequency is substantially lower than the single-core turbo frequency.
These observations explain quantitatively why measured efficiencies at n
= 8 cannot approach 100\% on this class of hardware, independently of
the algorithm.

\section{Code Implementation}

The project consists of four C programs sharing a single header
(common.h) that contains the dataset loader, the parameter estimator,
the simulation kernel and the reporting functions. Sharing the kernel
guarantees that all versions execute literally the same arithmetic. The
programs are built with GCC (MinGW) and
Microsoft MPI on Windows; all paths and the sequential reference time
are configurable from the command line.

\subsection{The Shared Kernel and the Seeding Scheme}

The kernel simulates the global trajectory indices \texttt{[sim\_begin, sim\_end)} of one
column and accumulates per-day sums. The seed line is the cornerstone of
the parallel design:

\begin{verbatim}
static void run_column_chunk(int c, int sim_begin, int sim_end,
    double start, double vol, double drift,
    double sum_d[DAYS], double sum2_d[DAYS]) {
  for (int sim = sim_begin; sim < sim_end; sim++) {
    unsigned int seed =
        (unsigned int)((long long)c * NUM_SIMS + sim + 1);
    /* ... 30 days x 1440 one-minute stochastic steps ... */
  }
}
\end{verbatim}

Because the seed depends only on the pair (column, global index), the
trajectory computed for a given (c, sim) is identical no matter which
version, rank or thread executes it, and any partition of the index
range yields partial sums that add exactly to the full-range sums.

\subsection{OpenMP Version --- Task Parallelism with Explicit Dependencies}

\begin{verbatim}
#pragma omp parallel num_threads(nthreads)
{
  #pragma omp single
  {
    for (int c = 0; c < NCOLS; c++) {
      #pragma omp task firstprivate(c) depend(out: params[c])
      { /* Stage A: estimate */
        estimate_params(data[c], nrows,
          &params[c][0], &params[c][1], &params[c][2]);
        free(data[c]); data[c] = NULL;
      }
      #pragma omp task firstprivate(c) \
        depend(in: params[c]) depend(out: sums[c])
      { /* Stage B: simulate */
        run_column_chunk(c, 0, NUM_SIMS, params[c][0],
          params[c][1], params[c][2], sums[c], sums2[c]);
      }
      #pragma omp task firstprivate(c) depend(in: sums[c])
      { /* Stage C: finalize */
        finalize_stats(sums[c], sums2[c], NUM_SIMS,
          all_mn[c], all_lo[c], all_hi[c]);
      }
    }
    /* fan-in: wait for all 21 tasks (7 cols x 3) */
    #pragma omp taskwait
  }
}
\end{verbatim}

Each task is the complete pipeline of one column --- one parallel path
of the dependency graph of Figure 3 --- and the depend clauses encode
the graph's edges directly in the source code. Tasks
write disjoint array slots, so no synchronization beyond the final
taskwait is needed.

\subsection{MPI Version --- Data Parallelism over Simulation Blocks}

\begin{verbatim}
int s0 = (int)((long long)NUM_SIMS * rank / size);
int s1 = (int)((long long)NUM_SIMS * (rank + 1) / size);
MPI_Bcast(params, NCOLS * 3, MPI_DOUBLE, 0, MPI_COMM_WORLD);
for (int c = 0; c < NCOLS; c++)
  run_column_chunk(c, s0, s1, params[c*3], params[c*3+1],
    params[c*3+2], local_sum[c], local_sum2[c]);
MPI_Reduce(local_sum, glob_sum, NCOLS*DAYS, MPI_DOUBLE,
  MPI_SUM, 0, MPI_COMM_WORLD);
MPI_Reduce(local_sum2, glob_sum2, NCOLS*DAYS,
  MPI_DOUBLE, MPI_SUM, 0, MPI_COMM_WORLD);
if (rank == 0) /* statistics finalized ONLY on reduced sums */
  for (int c = 0; c < NCOLS; c++)
    finalize_stats(glob_sum[c], glob_sum2[c],
                   NUM_SIMS, ...);
\end{verbatim}

Every rank processes its block of trajectories for all seven columns;
the only communication events are the initial broadcast and the final
sum-reduction of the raw accumulators.

\subsection{Hybrid Version}

The hybrid program composes the two mechanisms orthogonally: the MPI
block decomposition of Section 6.3 selects the local index range [s0,
s1), and an OpenMP task region executes the seven column pipelines over
that range inside each rank. Because the model parameters are estimated
once on rank 0 and broadcast to every rank (Section 3.1), the per-column
estimation stage is no longer part of the in-rank task graph; inside a
rank each column is therefore a two-stage pipeline instead of three: a
simulation task that produces the column's partial sums (depend(out:
sums[c])), followed by a local-reduce task that copies those sums
into the rank's MPI send buffer (depend(in: sums[c])). The
local-reduce task cannot start until the simulation task of the same
column has finished, because it reads sums[c] that the simulation
task wrote --- a real producer--consumer dependency; different columns
remain on independent paths, so up to seven column pipelines run
concurrently inside a rank, joined by a final taskwait (7 columns $\times$ 2 =
14 tasks per rank) before the reduction. MPI is initialized with
MPI\_Init\_thread(\ldots, MPI\_THREAD\_FUNNELED, \ldots) and all MPI
calls are performed by the master thread outside the parallel region, as
required by the funnelled threading level. The total worker count is
ranks $\times$ threads.

\subsection{Timing and Reporting}

All versions measure wall-clock time (QueryPerformanceCounter on
Windows; MPI\_Wtime and omp\_get\_wtime in the parallel versions), never
CPU time, so that the sequential baseline and the parallel measurements
are directly comparable. In the MPI versions the simulation time is
reduced with MPI\_MAX across ranks: the slowest rank defines T,
consistently with the fact that the reduction cannot complete earlier.
Every run appends one row (configuration, workers, T, speedup,
efficiency, throughput) to a shared CSV file from which the tables and
plots of Chapter 7 are generated.

\subsection{Correctness Verification}

The position-based seeding makes correctness directly testable: all
versions must produce identical statistics. This was verified on the
full benchmark workload: the eleven runs of Chapter 7 (sequential;
OpenMP with 2, 4, 7 threads; MPI with 2, 4, 8 ranks; hybrid 2$\times$2, 2$\times$4,
4$\times$2) report identical forecast values and identical confidence intervals
to all printed digits for every column and every day --- for example,
the Voltage day-30 forecast is 238.8127 V with 95\% interval
[229.1495, 248.4759] in every report. In addition, the additivity of
the block decomposition was unit-tested by comparing a full-range kernel
invocation with the sum of three uneven sub-range invocations: the
maximum absolute difference across all accumulators was exactly 0.0.

\section{Performance Evaluation}

\subsection{Setup and Methodology}

All experiments were executed on a single laptop with an Intel Core
i5-12450H CPU (12th generation, hybrid architecture: 4 performance cores
with two hardware threads each plus 4 efficiency cores, 8 physical cores
in total), 16 GB of RAM, running 64-bit Windows. The programs were
compiled with
GCC (MinGW-w64) at optimization level -O2; the MPI versions use
Microsoft MPI and are launched with mpiexec. The workload is fixed at
NUM\_SIMS = 50,000 trajectories per column, 7 columns, 30 days, 1,440
steps per day. The machine was connected to AC power with the Windows
power mode set to best performance and no foreground applications
running. Timings refer to the simulation phase measured with wall-clock
timers; for MPI runs the time of the slowest rank is reported. Anomalous
single measurements were re-executed and the best observed time is
reported; in particular the 2-rank MPI configuration was measured twice
(441.1 s and 408.1 s) and the better value is used, the residual gap
from the theoretical expectation being attributed to core placement
(Section 5.3).

\subsection{Metrics for Performance Evaluation}

Three standard metrics are used. The execution time T is the simulation
wall time with n processing elements. The speedup is $S = T_1 / T$, where $T_1 =
675.140$ s is the sequential baseline; linear (ideal) speedup corresponds
to $S = n$. The efficiency is $E = S / n$, the fraction of ideal speedup
actually achieved; $E = 100\%$ means perfect utilization of the n
processing elements. A throughput value (simulated trajectories per
second) is also recorded for convenience. For the task-parallel level a
theoretical bound complements the metrics: the makespan cannot be
smaller than the longest of the seven tasks, hence $S \leq T_1 / \max_c(t_c) \approx
4.56$ regardless of the thread count.

\subsection{Results and Analysis}

Tables 1--3 report the measured times, speedups and efficiencies of the
three required configurations; Figure 5 compares the execution times,
Figure 6 the speedups against the ideal line and Figure 7 the
efficiencies.

\emph{Table 1: Pure OpenMP (task parallelism). Baseline $T_1 = 675.140$ s.}

{\def\LTcaptype{none}
\begin{longtable}[]{@{}
  >{\centering\arraybackslash}p{(\linewidth - 6\tabcolsep) * \real{0.2500}}
  >{\centering\arraybackslash}p{(\linewidth - 6\tabcolsep) * \real{0.2500}}
  >{\centering\arraybackslash}p{(\linewidth - 6\tabcolsep) * \real{0.2500}}
  >{\centering\arraybackslash}p{(\linewidth - 6\tabcolsep) * \real{0.2499}}@{}}
\toprule\noalign{}
\textbf{Threads (n)} & \textbf{T(s)} & \textbf{Speedup S} & \textbf{Efficiency E} \\
\midrule\noalign{}
\endhead
\bottomrule\noalign{}
\endlastfoot
1 (sequential) & 675.140 & 1.000 & 100.0\% \\
2 & 480.664 & 1.405 & 70.2\% \\
4 & 351.293 & 1.922 & 48.0\% \\
7 & 254.328 & 2.655 & 37.9\% \\
\end{longtable}
}

\emph{Table 2: Pure MPI (data parallelism). Baseline $T_1 = 675.140$ s.}

{\def\LTcaptype{none}
\begin{longtable}[]{@{}
  >{\centering\arraybackslash}p{(\linewidth - 6\tabcolsep) * \real{0.2500}}
  >{\centering\arraybackslash}p{(\linewidth - 6\tabcolsep) * \real{0.2500}}
  >{\centering\arraybackslash}p{(\linewidth - 6\tabcolsep) * \real{0.2500}}
  >{\centering\arraybackslash}p{(\linewidth - 6\tabcolsep) * \real{0.2499}}@{}}
\toprule\noalign{}
\textbf{Ranks (n)} & \textbf{T(s)} & \textbf{Speedup S} & \textbf{Efficiency E} \\
\midrule\noalign{}
\endhead
\bottomrule\noalign{}
\endlastfoot
1 (sequential) & 675.140 & 1.000 & 100.0\% \\
2 & 354.248 & 1.906 & 95.3\% \\
4 & 250.278 & 2.698 & 67.4\% \\
8 & 185.957 & 3.631 & 45.4\% \\
\end{longtable}
}

\emph{Table 3: Hybrid MPI $\times$ OpenMP. Baseline $T_1 = 675.140$ s.}

{\def\LTcaptype{none}
\begin{longtable}[]{@{}
  >{\raggedright\arraybackslash}p{(\linewidth - 8\tabcolsep) * \real{0.2500}}
  >{\centering\arraybackslash}p{(\linewidth - 8\tabcolsep) * \real{0.1500}}
  >{\centering\arraybackslash}p{(\linewidth - 8\tabcolsep) * \real{0.2000}}
  >{\centering\arraybackslash}p{(\linewidth - 8\tabcolsep) * \real{0.2000}}
  >{\centering\arraybackslash}p{(\linewidth - 8\tabcolsep) * \real{0.2000}}@{}}
\toprule\noalign{}
\textbf{Configuration} & \textbf{Workers (n)} & \textbf{T(s)} & \textbf{Speedup S} & \textbf{Efficiency E} \\
\midrule\noalign{}
\endhead
\bottomrule\noalign{}
\endlastfoot
2 ranks $\times$ 2 threads & 4 & 277.709 & 2.431 & 60.8\% \\
2 ranks $\times$ 4 threads & 8 & 203.475 & 3.318 & 41.5\% \\
4 ranks $\times$ 2 threads & 8 & 201.589 & 3.349 & 41.9\% \\
\end{longtable}
}

\begin{quote}
\emph{Figure 5: Simulation time by configuration (NUM\_SIMS = 50,000,
Intel i5-12450H). The MPI 8-rank configuration is the fastest overall at
186.0 s.}
\end{quote}

\includegraphics[width=4.0625in,height=2.94167in]{media/image7.png}

\begin{quote}
\emph{Figure 6: Speedup versus number of processing elements, with the
ideal line S = n and the task-parallel upper bound T/max(task) = 4.56.
The MPI/hybrid configurations dominate at every worker count.}
\end{quote}

\includegraphics[width=4.0625in,height=2.94167in]{media/image8.png}

\begin{quote}
\emph{Figure 7: Efficiency versus number of processing elements.
Efficiency decreases with n on all curves, with a marked drop beyond
four workers when the efficiency cores of the CPU come into play.}
\end{quote}

The headline result is that the data-parallel MPI version scales best on
this machine: the 8-rank configuration completes in 186.0 s ($S = 3.63$, $E
= 45.4\%$), faster than the best pure OpenMP configuration (254.3 s with
7 threads) and the best hybrid configuration (201.6 s with 4 ranks $\times$ 2
threads). This reflects the strength of data parallelism: equal-sized
simulation blocks give near-perfect load balance and require only two
short collective operations, so MPI keeps improving as ranks are added
up to the core count, whereas the task-parallel level is capped at seven
unequal tasks.

The pure OpenMP curve follows the task-granularity theory
quantitatively. At 2 threads, efficiency is reasonable (70.2\%) because
the seven tasks pack into two workers. At 4 threads the speedup is 1.92:
packing seven unequal tasks (62--148 s, Figure 8) into four workers
requires two waves and the measured makespan matches the scheduling
bound almost exactly. At 7 threads the speedup is 2.66 against the
imbalance bound of 4.56; the gap is explained by the recorded per-task
times, which inflate uniformly to 250 s under full concurrency --- the
signature of the lower all-core turbo frequency and of tasks placed on
efficiency cores. The task level therefore hits three stacked ceilings:
the task count (7), the longest task (4.56$\times$) and the hardware
frequency/heterogeneity, in that order of severity.

The pure MPI curve shows the complementary behaviour. At 2 ranks the
speedup is 1.91 (95.3\% efficiency) --- almost perfect, because two
equal blocks map cleanly onto two performance cores with negligible
communication. At 4 ranks the speedup rises to 2.70 (67.4\%), and at 8
ranks to 3.63 (45.4\%): MPI
keeps improving as ranks are added, since the block decomposition stays
perfectly balanced. The efficiency drop beyond four workers is a
hardware effect --- four of the eight ranks must run on the slower
efficiency cores, and the all-core turbo frequency is lower --- not a
load-balancing or communication limit, which remain negligible
throughout.

Finally, the comparison between the two 8-worker hybrid configurations
isolates the value of rank placement: 4$\times$2 (201.6 s) slightly outperforms
2$\times$4 (203.5 s) with the same worker count, because four ranks anchor the
computation on the four performance cores and split the trajectory load
evenly among them, whereas with two ranks each process must spread its
four threads across heterogeneous cores and the task imbalance weighs
more inside each rank.

\subsection{Challenges and Limitations}

Three limitations qualify the results. First, the experiments run on a
single machine: the MPI level is exercised across processes, not across
physical nodes, so network costs are absent; on a real cluster the
constant 3.6 KB communication volume would still be negligible, but
latency would add a small fixed cost per run. Second, the CPU
heterogeneity (P-cores and E-cores) caps the achievable efficiency at
high worker counts independently of the software design; the reported
efficiencies at n = 8 measure the hardware as much as the program.
Third, single-run timing noise on a desktop operating system is
non-negligible; the methodology mitigates it by controlled conditions
and re-execution of anomalous points, but a production-grade study would
report averages with confidence intervals over many repetitions. A
natural extension of the study is a problem-size scaling analysis
(varying NUM\_SIMS) to show the amortization of the constant overheads.

\subsection{Conclusions}

The project implemented and compared four versions of a Monte Carlo
forecasting workload: sequential, task-parallel with OpenMP,
data-parallel with MPI, and hybrid. All versions provably compute
identical results thanks to a position-based random seeding scheme. The
experimental study on an 8-core hybrid CPU shows that the data-parallel
MPI level scales with near-perfect balance up to the number of
homogeneous cores, that the task-parallel OpenMP level is bounded by its
fixed granularity (7 tasks) and by task imbalance (upper bound 4.56$\times$),
and that the hybrid composition outperforms both pure approaches,
reaching a 3.63$\times$ speedup (186.0 s versus 675.1 s) with the 8-rank MPI
configuration, while the best hybrid (4$\times$2) reaches 3.35$\times$ (201.6 s).
Beyond the raw numbers, the measurements quantify three general lessons
of parallel computing on modern hardware: makespan is governed by the
slowest worker, task granularity bounds task-level speedup, and
process-to-core mapping on heterogeneous CPUs is a first-order
performance factor.

\section{References}

\begin{enumerate}
\item G. H\'{e}brail and A. B\'{e}rard, ``Individual household electric power
consumption Data Set'', UCI Machine Learning Repository, 2012.

\item Message Passing Interface Forum, ``MPI: A Message-Passing
Interface Standard'', \url{https://www.mpi-forum.org}.

\item OpenMP Architecture Review Board, ``OpenMP Application
Programming Interface Specification'', \url{https://www.openmp.org}.

\item P. Pacheco, ``An Introduction to Parallel Programming'', Morgan
Kaufmann, 2011.

\item J. S. Rosenthal, ``Parallel computing and Monte Carlo
algorithms'', Far East Journal of Theoretical Statistics, vol. 4, pp.
207--236, 2000.

\item I. Foster, ``Designing and Building Parallel Programs'',
Addison-Wesley, 1995.
\end{enumerate}

\end{document}